% mais de 100 paginas twoside
% menos de 100 paginas oneside
\documentclass[12pt, a4paper,oneside]{book}

% margens segundo abnt
\usepackage[top=3cm,left=3cm,right=2cm,bottom=2cm]{geometry}

% espaçamento entre linhas
\linespread{1.3}

% espacamento entre parágrafos
\setlength{\parindent}{1.25cm}

% colocar parágrafo no começo das seções
\usepackage{indentfirst}

% aceitar acentos, cedilhas, tils
\usepackage[utf8]{inputenc}
\usepackage[brazil]{babel}

% simbolos matemáticos
\usepackage{amsmath}
\usepackage{amsthm}
\usepackage{amsfonts}
\usepackage{amssymb}

% gerador de texto aleatório
\usepackage{lipsum}

% para subitem
\usepackage{enumitem}

% para sub sub item
\usepackage{outlines}

% para incluir figuras
\usepackage{graphicx}

% arrumar tabelas e figuras nos seus lugares
\usepackage{float}

% bibliografia e estilo
\usepackage[round]{natbib}
\bibliographystyle{apa}

% caminho das figuras
\graphicspath{{Figuras/}}

% insere paginas em pdf (usar para inserir a ficha) 
\usepackage{pdfpages}

% retira primeira página do capítulo
\usepackage{etoolbox}
\patchcmd{\chapter}{plain}{empty}{}{}
\patchcmd{\part}{plain}{empty}{}{}

% para criar teoremas
\newtheorem{thm}{Teorema}[section]
\newtheorem{exe}{Exemplo}[section]
\newtheorem{dfn}{Definição}[section]
\newtheorem{prob}{Problema}[section]
\newtheorem{cor}{Corolário}[section]
\newtheorem{prop}{Proposição}[section]
\newtheorem{lem}{Lema} [section]
\newcounter{contar}

% tira a palavra capitulo dos capitulos (ABNT)
\usepackage{titlesec}
\titleformat{\chapter}{\huge\bf}{\thechapter}{20pt}{\huge\bf}


\begin{document}
% numera primeiras páginas em romano
\frontmatter

%---------- CAPA -------------

% Não numera a capa
\thispagestyle{empty}

\begin{center}
\sc{\large{Universidade Federal do Triângulo Mineiro - UFTM} } \\
\vspace{0.2cm}
\sc{\includegraphics[scale=0.35]{logouftm.jpg} } \\
\vspace{0.5cm}
\sc{\small{Mestrado Profissional em Matemática em Rede Nacional - PROFMAT}} \\
\vspace{0.2cm}
\sc{\includegraphics[scale=1.2]{logoProfmat.jpg} } \\
\vspace{0.5cm}
\sc{\small{Dissertação de Mestrado}} \\

\vspace{3cm}

\sc{\Large{Coloque aqui o título da sua dissertação}}

\vspace{2.5cm}

\sc{\Large{Coloque aqui seu nome}}

\vspace{3.5cm}

\textbf{Uberaba - Minas Gerais} \\
Julho de 2016
\end{center}


%---------- FOLHA DE ROSTO -------------
\newpage
\thispagestyle{empty}

\begin{center}
\sc{\Large{Coloque aqui o título da sua dissertação}}

\vspace{4cm}

\large{Coloque aqui seu nome}
\end{center}

\vspace{4cm}

\begin{flushright}
\begin{minipage}{9.1cm}
Dissertação de Mestrado apresentada à \mbox{Comissão} 
Acadêmica Institucional do PROFMAT-UFTM como 
requisito parcial para obtenção 
do título de Mestre em \mbox{Matemática}.

\vspace{0.5cm}
\textbf{Orientador}: Prof. Dr. Orientador Super Legal.

\end{minipage}
\end{flushright}
 
\vspace{8cm}


\begin{center}
\textbf{Uberaba - Minas Gerais} \\
Agosto de 2016
\end{center}

%-------------Ficha----------------

\includepdf{Ficha/fichaCatalografica.pdf}

%---------- BANCA EXAMINADORA -------------
\includepdf{Assinaturas/Assinaturas.pdf}

%-----------Dedicatória----------------
\include{dedicatoria}

%------------Agradecimentos------------
\include{agradecimentos}

%------------Citação-------------------
\include{citacao}


%--------------Resumo-------------------
\pagestyle{myheadings} % tira texto do cabeçalho
\newpage
\thispagestyle{empty}
\chapter*{Resumo}

Lorem ipsum dolor sit amet, consectetur adipiscing elit. Duis eu ligula neque. Sed vehicula enim ac neque varius, vitae congue magna pretium. Fusce tristique, arcu in tempus auctor, ligula justo malesuada justo, at fermentum arcu dui nec quam. Aliquam elit metus, euismod quis justo eu, sodales suscipit elit. Nulla tempor odio blandit tincidunt feugiat. Nulla quis sodales mi. Praesent ac accumsan sapien. Ut sollicitudin nisl sit amet mollis ultrices. Ut tincidunt turpis id interdum rutrum. Vivamus a nisl vel dolor molestie ullamcorper.
 
\noindent  {\bf Palavras-chave:} Lorem, Ipsum, Loremipsum, Ipsumlorem.


%-------------Abstract------------------
\newpage
\thispagestyle{empty}
\chapter*{Abstract}

Lorem ipsum dolor sit amet, consectetur adipiscing elit. Duis eu ligula neque. Sed vehicula enim ac neque varius, vitae congue magna pretium. Fusce tristique, arcu in tempus auctor, ligula justo malesuada justo, at fermentum arcu dui nec quam. Aliquam elit metus, euismod quis justo eu, sodales suscipit elit. Nulla tempor odio blandit tincidunt feugiat. Nulla quis sodales mi. Praesent ac accumsan sapien. Ut sollicitudin nisl sit amet mollis ultrices. Ut tincidunt turpis id interdum rutrum. Vivamus a nisl vel dolor molestie ullamcorper. 

\noindent  {\bf Keywords:} Lorem, Ipsum, Loremipsum, Ipsumlorem.


%-------------Índice--------------------
\newpage
\thispagestyle{empty}
\tableofcontents


%-------------Lista de figuras----------
\newpage
\listoffigures


%-------------Lista de Tabelas----------
\newpage
\listoftables

%-------------Introdução----------------
\mainmatter
\include{introducao}

%-------------Capítulo 1-----------------
\include{capitulo1}

%-------------Capítulo 2-----------------
\include{capitulo2}

%-------------Conclusão-----------------
\include{conclusao}

\nocite{khan,da2005aplicaccao,gama2004,lior2014,dados1}

%-------------Bibliografia------------------
\newpage
\renewcommand{\refname}{REFERÊNCIAS BIBLIOGRÁFICAS}
\addcontentsline{toc}{chapter}{REFERÊNCIAS BIBLIOGRÁFICAS}
\bibliography{bibliografia}

%-----------APENDICES---------------------
\renewcommand{\refname}{APÊNDICES}
\addcontentsline{toc}{chapter}{APÊNDICES}
\appendix % Cue to tell LaTeX that the following "chapters" are Appendices

\include{Apendices/apendiceA}
\include{Apendices/apendiceB}

\end{document}